\chapter{Complexity — Oblivious Candidate}
%%% generated by the blueprint pipeline from TCSlib.Complexity.TuringMachine.Robustness.ObliviousCandidate
\section{Overview}

This module assembles the concrete candidate machine for the oblivious-simulation
theorem (Arora--Barak, Remark~1.7 and Exercise~1.5): a payload alphabet for
head-centered virtual tapes, the data transducer that replays one step of a source
machine per sweep of a fixed physical schedule, and the binary simulator obtained by
decorating the length-only schedule with that transducer.  It also proves the exact
run identities consumed by the later cost and correctness analysis: the binary budget
counter and its borrow sweeps, a general one-lane transduction lemma, the
head-centered payload correspondence for virtual heads, the captured-clock stage, and
the guide-tape macrostep cycle with its exact duration of $6R+8$ transitions.

\section{Declarations}

\begin{definition}[Payload alphabet for virtual cells]
\lean{Complexity.OblPayload}
\label{Complexity.OblPayload}
\leanok
A payload is a pair consisting of an optional bit (the represented cell content,
absent for a blank cell) and a tag from a three-element set: tag $0$ marks an
ordinary cell, tag $1$ the left virtual input boundary, and tag $2$ the right
virtual input boundary.  Payloads describe the cells of virtual tapes whose
logical head sits at a fixed marked origin, so that moving a virtual head shifts
the represented tape while the physical schedule stays fixed.
\end{definition}

\begin{definition}[Blank payload]
\lean{Complexity.blankPayload}
\label{Complexity.blankPayload}
\leanok
\uses{Complexity.OblPayload}
There is a distinguished payload consisting of a blank cell content together with the
ordinary tag $0$; it also serves as the neighbor value just outside a sweep.
\end{definition}

\begin{definition}[Decoding of a physical cell]
\lean{Complexity.dataCell}
\label{Complexity.dataCell}
\leanok
\uses{Complexity.OblPayload, Complexity.OblSymbol, Complexity.blankPayload}
Given an optional symbol of the extended simulation alphabet, its decoding is a pair
of payloads: a data-cell symbol carrying a current payload and a cached left-neighbor
payload decodes to exactly that pair, while every other symbol, including the untouched
physical blank, decodes to the pair of blank payloads.
\end{definition}

\begin{definition}[Finite data state of the simulator]
\lean{Complexity.OblData}
\label{Complexity.OblData}
\leanok
\uses{Complexity.OblPayload, Turing.FinTM}
Let $M$ be a binary multi-tape Turing machine with $k$ work tapes.  The associated
finite data state consists of an optional control state of $M$ (absent once $M$ has
halted), a saved one-bit answer, and two vectors of $k+1$ payloads: the reads captured
at macrostep entry and the neighbor registers used during a sweep.  All tape
coordinates remain on tapes, so this state is finite.
\end{definition}

\begin{definition}[Source action selected from payload reads]
\lean{Complexity.obliviousSourceAction}
\label{Complexity.obliviousSourceAction}
\leanok
\uses{Complexity.OblPayload, Turing.Action, Turing.FinTM}
Let $M$ be a binary machine with $k$ work tapes, let $q$ be an optional control state
of $M$, and let a vector of $k+1$ payloads be given (one per work tape and one for the
virtual input).  If $q$ is absent (the machine has halted), the selected action is the
identity: no writes, no output, no head movement, and a halted successor state.  If
$q$ is a live state, the selected action is the transition of $M$ from $q$ on the bit
carried by the virtual-input payload and the bits carried by the work-tape payloads.
\end{definition}

\begin{definition}[Virtual head directions of a macrostep]
\lean{Complexity.obliviousSourceMove}
\label{Complexity.obliviousSourceMove}
\leanok
\uses{Complexity.OblPayload, Complexity.obliviousSourceAction, Turing.FinTM}
Let $M$ be a binary machine with $k$ work tapes, let $q$ be an optional control state,
and let a vector of $k+1$ payload reads be given.  For each work tape the virtual head
direction is the movement prescribed by the action selected from $q$ and the reads;
the virtual input head direction is that action's input movement, clamped to
stationary when the input payload carries the left-boundary tag and the movement is
leftward, or the right-boundary tag and the movement is rightward.  These directions
select shifted payloads and never move a physical head.
\end{definition}

\begin{definition}[Initial data state]
\lean{Complexity.obliviousDataInit}
\label{Complexity.obliviousDataInit}
\leanok
\uses{Complexity.OblData, Complexity.blankPayload, Turing.FinTM}
For a binary machine $M$, the initial data state consists of the live initial control
state of $M$, the saved answer bit set to false, and blank payloads in both register
vectors.
\end{definition}

\begin{definition}[Data transducer along the schedule]
\lean{Complexity.obliviousVisit}
\label{Complexity.obliviousVisit}
\leanok
\uses{Complexity.OblData, Complexity.OblPayload, Complexity.OblPhase, Complexity.OblSymbol, Complexity.blankPayload, Complexity.clockBit, Complexity.dataCell, Complexity.obliviousSourceAction, Complexity.obliviousSourceMove, Turing.FinTM}
Let $W$ and $M$ be binary machines and $a$ a padding constant.  Given the current
phase of the schedule controller, a data state for $M$, the symbol under the input
head, the reads of the controller's tapes, and the reads of the $k+1$ data tapes, this
transducer returns an updated data state, an optional write for each data tape, and an
optional output symbol.  During the copy phases the input bit is written with explicit
boundary payloads; at a marked origin with a positive counter test the source
transition is selected and its writes performed; the forward sweep caches left
neighbors; the backward sweep shifts each virtual tape according to its selected
direction; the return to the origin installs the source successor state (possibly
halted); and only the final failed unary-counter test emits the stored answer bit.
\end{definition}

\begin{definition}[The concrete oblivious simulator]
\lean{Complexity.obliviousCandidate}
\label{Complexity.obliviousCandidate}
\leanok
\uses{Complexity.decorateTM, Complexity.oblEmbed, Complexity.obliviousDataInit, Complexity.obliviousSchedule, Complexity.obliviousVisit, Complexity.parallelTM, Turing.FinTM}
Given binary machines $W$ (the clock witness) and $M$ (the source) and a padding
constant $a$, the candidate simulator is obtained by decorating the length-only
schedule controller built from $W$ and $a$ with $k+1$ additional data tapes driven by
the data transducer, starting from the initial data state, and then realizing the
finite alphabet over the binary alphabet by simultaneous transverse binary blocks;
this coding retains the exact schedule.
\end{definition}

\begin{lemma}[Obliviousness of the candidate]
\lean{Complexity.obliviousCandidate_oblivious}
\label{Complexity.obliviousCandidate_oblivious}
\leanok
\uses{Complexity.decorateTM_oblivious, Complexity.obliviousCandidate, Complexity.obliviousSchedule_oblivious, Complexity.obliviousSchedule_output, Complexity.parallelTM_oblivious, Turing.FinTM, Turing.FinTM.Oblivious}
\difficulty{3}
For all binary machines $W$ and $M$ and every padding constant $a$, the candidate
simulator built from $W$, $M$, and $a$ is oblivious: its head positions at every time
step depend only on the length of the input.  This holds on all binary inputs of every
length, including the empty input, and uses no assumption about the computations of
$W$ or $M$.
\end{lemma}

\begin{lemma}[Idle source action]
\lean{Complexity.obliviousSourceAction_idle}
\label{Complexity.obliviousSourceAction_idle}
\leanok
\uses{Complexity.OblPayload, Complexity.obliviousSourceAction, Turing.FinTM}
\difficulty{1}
For every binary machine $M$ and every vector of payload reads, the action selected in
the halted state is the identity action: it performs no writes, emits no output, moves
no head, and keeps the halted successor state.
\end{lemma}

\begin{lemma}[Idle virtual heads]
\lean{Complexity.obliviousSourceMove_idle}
\label{Complexity.obliviousSourceMove_idle}
\leanok
\uses{Complexity.OblPayload, Complexity.obliviousSourceAction, Complexity.obliviousSourceMove, Turing.FinTM}
\difficulty{3}
For every binary machine $M$ and every vector of payload reads, each of the $k+1$
virtual head directions selected in the halted state is stationary.
\end{lemma}

\begin{definition}[Little-endian value of a budget word]
\lean{Complexity.budgetValue}
\label{Complexity.budgetValue}
\leanok
Given a list of bits, its little-endian numeric value is defined recursively: the
empty word has value $0$, and a word with first bit $b$ and remaining word of value
$v$ has value $2v + b$.  Zero-padded words are permitted and high zero bits do not
change the value.
\end{definition}

\begin{lemma}[Value of the canonical bit representation]
\lean{Complexity.budgetValue_bits}
\label{Complexity.budgetValue_bits}
\leanok
\uses{Complexity.budgetValue}
\difficulty{4}
For every natural number $n$, the little-endian value of the canonical binary
representation of $n$ equals $n$.
\end{lemma}

\begin{definition}[Fixed-width borrow pass]
\lean{Complexity.budgetBorrow}
\label{Complexity.budgetBorrow}
\leanok
Given an incoming borrow flag and a bit word, the borrow pass processes every bit of
the word: a bit $b$ met with borrow $c$ is rewritten to $b \oplus c$ and the borrow
becomes $c \wedge \lnot b$.  It returns the final borrow flag together with the
rewritten word, always retaining the full width of the word even when leading high
bits become zero.
\end{definition}

\begin{lemma}[Cleared borrow is the identity]
\lean{Complexity.budgetBorrow_false}
\label{Complexity.budgetBorrow_false}
\leanok
\uses{Complexity.budgetBorrow}
\difficulty{3}
For every bit word, the borrow pass entered with a cleared borrow flag returns a
cleared flag and leaves every bit of the word unchanged.
\end{lemma}

\begin{lemma}[Borrow pass preserves width]
\lean{Complexity.budgetBorrow_length}
\label{Complexity.budgetBorrow_length}
\leanok
\uses{Complexity.budgetBorrow}
\difficulty{3}
For every borrow flag and every bit word, the word produced by the borrow pass has
exactly the length of the input word.
\end{lemma}

\begin{lemma}[Underflow exactly at zero]
\lean{Complexity.budgetBorrow_underflow}
\label{Complexity.budgetBorrow_underflow}
\leanok
\uses{Complexity.budgetBorrow, Complexity.budgetBorrow_false, Complexity.budgetValue}
\difficulty{4}
For every bit word $w$, the borrow pass entered with a set borrow flag ends with a set
flag if and only if the little-endian value of $w$ is $0$; this holds also for words
padded with high zero bits.
\end{lemma}

\begin{lemma}[Successful borrow subtracts one]
\lean{Complexity.budgetBorrow_value}
\label{Complexity.budgetBorrow_value}
\leanok
\uses{Complexity.budgetBorrow, Complexity.budgetBorrow_false, Complexity.budgetValue}
\difficulty{4}
For every bit word $w$ whose little-endian value is positive, the word produced by the
borrow pass entered with a set flag has little-endian value exactly one less than the
value of $w$.
\end{lemma}

\begin{definition}[Single-lane configuration]
\lean{Complexity.laneCfg}
\label{Complexity.laneCfg}
\leanok
\uses{Turing.Cfg, Turing.FinTM.sweepTape}
Let a base configuration of a $k$-tape machine be given, together with a distinguished
work tape (the lane), an optional control state, an integer coordinate $z$, and two
lists of optional symbols: the already-processed cells in reverse order and the
unprocessed cells.  The lane configuration agrees with the base on every other tape,
on the input head, and on the output, while the lane carries the zipper tape
determined by $z$ and the two lists and its head sits at $z$.
\end{definition}

\begin{definition}[Single-lane action]
\lean{Complexity.laneAction}
\label{Complexity.laneAction}
\leanok
\uses{Turing.Action}
Given a designated work tape, a control state, an optional symbol, and a direction,
the single-lane action writes that symbol and moves in that direction on the
designated tape only, enters the given state, keeps the input head stationary, leaves
every other work tape untouched, and emits no output.
\end{definition}

\begin{lemma}[Lane read]
\lean{Complexity.laneCfg_read}
\label{Complexity.laneCfg_read}
\leanok
\uses{Complexity.laneCfg, Turing.Cfg, Turing.Cfg.workTapeSymbols, Turing.FinTM.sweepTape_read}
\difficulty{1}
In every single-lane configuration, the symbol read on the active lane is the first
entry of the unprocessed list, or blank if that list is empty.
\end{lemma}

\begin{lemma}[Right-moving lane step]
\lean{Complexity.laneCfg_right}
\label{Complexity.laneCfg_right}
\leanok
\uses{Complexity.laneAction, Complexity.laneCfg, Turing.Action.apply, Turing.Cfg, Turing.FinTM.sweepTape_right, Turing.moveInputPos_zero}
\difficulty{4}
Let a single-lane configuration be given whose unprocessed list begins with an entry
$s$.  Applying the single-lane action that writes an entry $s'$ on the lane, moves
right, and enters a state $q$ yields again a single-lane configuration: the state
becomes $q$, the coordinate advances by one, $s'$ is pushed onto the processed list,
$s$ is removed from the unprocessed list, and all other tapes and heads are unchanged.
\end{lemma}

\begin{lemma}[Exact cost of a one-lane transduction]
\lean{Complexity.lane_run}
\label{Complexity.lane_run}
\leanok
\uses{Complexity.budgetVisit, Complexity.laneAction, Complexity.laneCfg, Complexity.laneCfg_read, Complexity.laneCfg_right, Turing.Action.apply, Turing.Cfg, Turing.Cfg.workTapeSymbols, Turing.FinTM.sweepFold, Turing.MultiTapeTM, Turing.MultiTapeTM.runFrom, Turing.MultiTapeTM.runFrom_succ_eq_step, Turing.MultiTapeTM.runFrom_zero}
\difficulty{5}
Let a $k$-tape machine, a designated lane, an encoding of transducer states into
machine states, an encoding of transducer letters into tape symbols, and a local
transducer step be given.  Suppose that from every encoded transducer state, whenever
the lane reads an encoded letter, the machine performs exactly the single-lane action
that writes the encoded output letter of the local step, moves right, and enters the
encoded updated transducer state.  Then running the machine for exactly $\ell$ steps
from a single-lane configuration whose unprocessed list begins with an encoded word of
length $\ell$ yields the single-lane configuration in which the transducer has been
folded over the whole word: the state is the encoded final transducer state, the
coordinate has advanced by $\ell$, the encoded output word has been moved in reverse
onto the processed list, and no inactive tape has changed.
\end{lemma}

\begin{definition}[One-bit borrow transducer]
\lean{Complexity.budgetVisit}
\label{Complexity.budgetVisit}
\leanok
Given a carry bit $c$ and a data bit $b$, the local borrow transducer produces the new
carry $c \wedge \lnot b$ and the output bit $b \oplus c$.
\end{definition}

\begin{lemma}[Folding the borrow transducer]
\lean{Complexity.budgetFold}
\label{Complexity.budgetFold}
\leanok
\uses{Complexity.budgetBorrow, Complexity.budgetVisit, Turing.FinTM.sweepFold}
\difficulty{3}
For every carry flag and every bit word, folding the one-bit borrow transducer over
the word is precisely the fixed-width borrow pass on that word.
\end{lemma}

\begin{lemma}[Budget-lane focusing of a schedule action]
\lean{Complexity.oblAction_budget}
\label{Complexity.oblAction_budget}
\leanok
\uses{Complexity.OblSymbol, Complexity.laneAction, Complexity.oblAction}
\difficulty{4}
Consider a machine with $k$ leading work tapes followed by three administrative tapes,
of which the first is the budget tape.  For every state $q$, optional symbol $s$, and
direction $d$, the schedule action that enters $q$, keeps the input head fixed, writes
$s$ with direction $d$ on the budget tape, and leaves the two remaining administrative
tapes and all $k$ leading tapes untouched equals the single-lane action on the budget
tape writing $s$, moving in direction $d$, and entering $q$.
\end{lemma}

\begin{lemma}[Borrow transition of the schedule controller]
\lean{Complexity.obliviousSchedule_borrow}
\label{Complexity.obliviousSchedule_borrow}
\leanok
\uses{Complexity.OblSymbol, Complexity.laneAction, Complexity.oblAction_budget, Complexity.obliviousSchedule, Turing.FinTM}
\difficulty{2}
Let $W$ be a binary machine and $a$ a padding constant, and consider the length-only
schedule controller built from them, whose first administrative tape is the budget
tape.  If the budget tape currently reads a bit symbol $b$, then the transition out of
the borrow control phase with carry $c$ is the single-lane action on the budget tape
that writes the bit $b \oplus c$, moves right, and enters the borrow phase with carry
$c \wedge \lnot b$.
\end{lemma}

\begin{lemma}[Exact run of the implemented borrow pass]
\lean{Complexity.obliviousSchedule_borrow_run}
\label{Complexity.obliviousSchedule_borrow_run}
\leanok
\uses{Complexity.OblPhase, Complexity.OblSymbol, Complexity.budgetBorrow, Complexity.budgetFold, Complexity.budgetVisit, Complexity.inputPayload, Complexity.laneCfg, Complexity.lane_run, Complexity.obliviousSchedule, Complexity.obliviousSchedule_borrow, Turing.Cfg, Turing.FinTM, Turing.MultiTapeTM.runFrom}
\difficulty{2}
Let $W$ be a binary machine, $a$ a padding constant, and let a bit word $w$ be laid
out as bit symbols on the budget lane of the schedule controller, starting in the
borrow phase with carry $c$.  Running the controller for exactly $|w|$ steps carries
out the fixed-width borrow pass on $w$: the phase becomes borrow with the resulting
borrow flag, the rewritten word is moved in reverse into the processed zone, the head
has advanced by $|w|$, and every other tape content and head position, including
copied input data in later decorations, is unchanged.
\end{lemma}

\begin{definition}[Virtual input payload]
\lean{Complexity.inputPayload}
\label{Complexity.inputPayload}
\leanok
\uses{Complexity.OblPayload, Turing.FinTM.bufferTape}
For a binary word $x$ and an integer coordinate $z$, the virtual input payload at $z$
consists of the buffered input symbol of $x$ at $z$ together with the tag $1$ at
$z = -1$ (the left boundary), the tag $2$ at $z = |x|$ (the right boundary), and the
ordinary tag $0$ elsewhere.
\end{definition}

\begin{definition}[Boundary tag at the input head]
\lean{Complexity.inputTag}
\label{Complexity.inputTag}
\leanok
\uses{Turing.Cfg}
For a configuration of a binary machine on an input $x$, the boundary tag at the
native input head is $1$ when the head sits at position $0$ (the left end), $2$ when
it sits at position $|x| + 1$ (the right end), and $0$ otherwise.
\end{definition}

\begin{lemma}[Payload at the virtual input head]
\lean{Complexity.inputPayload_head}
\label{Complexity.inputPayload_head}
\leanok
\uses{Complexity.inputPayload, Complexity.inputTag, Turing.Cfg, Turing.Cfg.inputSymbol, Turing.FinTM.bufferTape_inputSymbol}
\difficulty{3}
For every configuration of a binary machine on an input $x$ with input head at
position $p$, the virtual input payload of $x$ at coordinate $p - 1$ is exactly the
pair of the symbol read by the native input head and the boundary tag of that
configuration; on the empty input both boundary tags are produced correctly.
\end{lemma}

\begin{definition}[Clipped virtual direction]
\lean{Complexity.clippedMove}
\label{Complexity.clippedMove}
\leanok
For an input head position $p \in \{0, 1, \dots, n+1\}$ and a direction $d$, the
clipped direction is stationary when $p = 0$ and $d$ is leftward, or when $p = n + 1$
and $d$ is rightward, and equals $d$ in every other case.  It expresses native input
clamping as a virtual tape-shift direction.
\end{definition}

\begin{lemma}[Clipping matches the native head displacement]
\lean{Complexity.clippedMove_correct}
\label{Complexity.clippedMove_correct}
\leanok
\uses{Complexity.clippedMove, Turing.FinTM.moveInputPos_neg_val, Turing.moveInputPos, Turing.moveInputPos_pos_of_ne_right}
\difficulty{4}
For every input head position $p \in \{0, 1, \dots, n+1\}$ and every direction $d$,
the position reached by the native clamped input move of $p$ in direction $d$, shifted
down by one, equals $p - 1$ plus the clipped direction viewed as an integer.
\end{lemma}

\begin{definition}[Head-centered source payloads]
\lean{Complexity.sourcePayload}
\label{Complexity.sourcePayload}
\leanok
\uses{Complexity.OblPayload, Complexity.inputPayload, Turing.Cfg}
Let a configuration of a binary machine with $k$ work tapes on an input $x$ be given,
and let $z$ be an integer offset.  The head-centered payload vector at $z$ assigns to
each work tape the symbol of that tape at signed distance $z$ from its head, with
ordinary tag $0$, and assigns to the final coordinate the virtual input payload of $x$
at the input head position minus one plus $z$; thus the last virtual tape is the
copied input and the earlier ones are the source work tapes.
\end{definition}

\begin{definition}[Source reads at the origin]
\lean{Complexity.sourceReads}
\label{Complexity.sourceReads}
\leanok
\uses{Complexity.OblPayload, Complexity.inputTag, Turing.Cfg, Turing.Cfg.inputSymbol, Turing.Cfg.workTapeSymbols}
For a configuration of a binary machine with $k$ work tapes, the read vector at a
macrostep origin assigns to each work tape the symbol under its head with ordinary tag
$0$, and to the final coordinate the symbol read by the input head paired with the
boundary tag of the configuration.
\end{definition}

\begin{lemma}[Origin payloads are the source reads]
\lean{Complexity.sourcePayload_origin}
\label{Complexity.sourcePayload_origin}
\leanok
\uses{Complexity.inputPayload_head, Complexity.sourcePayload, Complexity.sourceReads, Turing.Cfg, Turing.Cfg.workTapeSymbols}
\difficulty{3}
For every configuration of a binary machine, the head-centered payload vector at
offset $0$ equals the read vector at the macrostep origin; reading the marked origin
therefore obtains all source reads simultaneously.
\end{lemma}

\begin{definition}[Totalized source action]
\lean{Complexity.sourceTotalAction}
\label{Complexity.sourceTotalAction}
\leanok
\uses{Turing.Action, Turing.Cfg, Turing.Cfg.inputSymbol, Turing.Cfg.workTapeSymbols, Turing.FinTM}
For a binary machine $M$ and a configuration $c$ of $M$, the totalized action is the
identity action (no writes, no output, no head movement, halted successor) when $c$ is
halted, and the transition action of $M$ on the current input and work-tape reads of
$c$ otherwise; halting is thereby interpreted as an idle macrostep.
\end{definition}

\begin{lemma}[The controller selects the source transition]
\lean{Complexity.obliviousSourceAction_correct}
\label{Complexity.obliviousSourceAction_correct}
\leanok
\uses{Complexity.obliviousSourceAction, Complexity.sourcePayload, Complexity.sourcePayload_origin, Complexity.sourceReads, Complexity.sourceTotalAction, Turing.Cfg, Turing.FinTM}
\difficulty{3}
For every binary machine $M$ and every configuration $c$ of $M$ on an input $x$, the
action selected by the data controller from the control state of $c$ and the
head-centered payloads at offset $0$ equals the totalized action of $M$ at $c$.
\end{lemma}

\begin{lemma}[Applying the totalized action is one step]
\lean{Complexity.sourceTotalAction_apply}
\label{Complexity.sourceTotalAction_apply}
\leanok
\uses{Complexity.sourceTotalAction, Turing.Action.apply, Turing.Cfg, Turing.FinTM, Turing.MultiTapeTM.step}
\difficulty{4}
For every binary machine $M$ and every configuration $c$ of $M$, applying the
totalized action of $M$ at $c$ to $c$ yields exactly one machine step of $M$ from $c$,
including the idle step performed after $M$ has halted.
\end{lemma}

\begin{definition}[Shift directions of an action]
\lean{Complexity.sourceShift}
\label{Complexity.sourceShift}
\leanok
\uses{Complexity.clippedMove, Turing.Action, Turing.Cfg}
Given a configuration of a machine with $k$ work tapes and an action, the shift vector
assigns to each work tape the direction prescribed by the action for that tape, and to
the final coordinate the clipped direction of the action's input movement at the
current input head position.
\end{definition}

\begin{lemma}[The controller selects the correct virtual shifts]
\lean{Complexity.obliviousSourceMove_correct}
\label{Complexity.obliviousSourceMove_correct}
\leanok
\uses{Complexity.clippedMove, Complexity.inputTag, Complexity.obliviousSourceAction_correct, Complexity.obliviousSourceMove, Complexity.sourcePayload, Complexity.sourcePayload_origin, Complexity.sourceReads, Complexity.sourceShift, Complexity.sourceTotalAction, Turing.Cfg, Turing.FinTM}
\difficulty{4}
For every binary machine $M$ and every configuration $c$ of $M$ on an input $x$, the
vector of virtual head directions selected by the controller from the control state of
$c$ and the head-centered payloads at offset $0$ equals the shift vector of the
totalized action of $M$ at $c$; in particular native input clamping and source work
head movement are reproduced exactly.
\end{lemma}

\begin{definition}[Payloads after the origin write]
\lean{Complexity.sourceWrittenPayload}
\label{Complexity.sourceWrittenPayload}
\leanok
\uses{Complexity.OblPayload, Complexity.inputPayload, Turing.Action, Turing.Cfg, Turing.Cfg.workTapeSymbols}
Given a configuration of a machine with $k$ work tapes, an action, and an integer
offset $z$, the written payload vector describes the tapes after performing the
action's writes at the old origin but before any shifting: at offset $0$ each
work-tape entry is the symbol written by the action (or the symbol under the head when
the action writes nothing), at every other offset it is the untouched tape content,
and the input entries are read-only, always given by the virtual input payload.
\end{definition}

\begin{lemma}[Head-centered simulation identity]
\lean{Complexity.sourcePayload_apply}
\label{Complexity.sourcePayload_apply}
\leanok
\uses{Complexity.clippedMove_correct, Complexity.sourcePayload, Complexity.sourceShift, Complexity.sourceWrittenPayload, Turing.Action, Turing.Action.apply, Turing.Cfg, Turing.Cfg.workTapeSymbols}
\difficulty{5}
Let $c$ be a configuration of a binary machine with $k$ work tapes, let $\alpha$ be an
action, let $z$ be an integer offset, and let $i$ be one of the $k+1$ virtual tape
coordinates.  Then the head-centered payload of the configuration obtained by applying
$\alpha$ to $c$, at offset $z$ and coordinate $i$, equals the written payload of $c$
and $\alpha$ at offset $z$ plus the shift direction of coordinate $i$: writing at the
old origin and then shifting each virtual tape by its displacement produces exactly
the payloads of the successor configuration, including stationary moves, boundary
attempts, and blank writes.
\end{lemma}

\begin{definition}[Parallel payload row]
\lean{Complexity.payloadRow}
\label{Complexity.payloadRow}
\leanok
\uses{Complexity.OblPayload, Complexity.blankPayload}
Given a family of payloads indexed by an integer coordinate and a tape index, a
coordinate $z$, and a caching flag, the row at $z$ assigns to each tape the pair of
its payload at $z$ together with, when the flag is set, its payload at $z - 1$ as
cached left neighbor, and the blank payload otherwise.
\end{definition}

\begin{definition}[Forward row transducer]
\lean{Complexity.payloadForward}
\label{Complexity.payloadForward}
\leanok
\uses{Complexity.OblPayload}
Given a left-neighbor register (one payload per tape) and a row, the forward
transducer outputs the row's payloads as the new register and rewrites the row so that
the incoming register is installed as its cached left neighbor.
\end{definition}

\begin{definition}[Backward row transducer]
\lean{Complexity.payloadBackward}
\label{Complexity.payloadBackward}
\leanok
\uses{Complexity.OblPayload, Complexity.blankPayload}
Given per-tape shift directions, a right-neighbor register, and a row, the backward
transducer outputs the row's payloads as the new register and rewrites each tape's
entry to its cached left neighbor, its current payload, or the incoming right register
according as that tape's direction is leftward, stationary, or rightward, clearing the
cache to blank; all physical heads move left while it runs.
\end{definition}

\begin{lemma}[One forward row]
\lean{Complexity.payloadForward_row}
\label{Complexity.payloadForward_row}
\leanok
\uses{Complexity.OblPayload, Complexity.payloadForward, Complexity.payloadRow}
\difficulty{1}
For every payload family and every coordinate $z$, applying the forward transducer
with register the payloads at $z - 1$ to the uncached row at $z$ produces the payloads
at $z$ as the new register together with the cached row at $z$.
\end{lemma}

\begin{lemma}[One backward row]
\lean{Complexity.payloadBackward_row}
\label{Complexity.payloadBackward_row}
\leanok
\uses{Complexity.OblPayload, Complexity.payloadBackward, Complexity.payloadRow}
\difficulty{3}
For every payload family $f$, every vector of directions $d$, and every coordinate
$z$, applying the backward transducer with directions $d$ and register the payloads at
$z + 1$ to the cached row at $z$ produces the payloads at $z$ as the new register
together with the uncached row at $z$ of the shifted family sending coordinate $z$ and
tape $i$ to the value of $f$ at $z$ plus the direction of $i$.
\end{lemma}

\begin{definition}[Payload zone]
\lean{Complexity.payloadZone}
\label{Complexity.payloadZone}
\leanok
\uses{Complexity.OblPayload, Complexity.payloadRow}
Given a payload family, a starting coordinate $z$, a caching flag, and a length $n$,
the zone is the list of the $n$ consecutive rows of the layout at coordinates
$z, z+1, \dots, z+n-1$, each cached or uncached according to the flag.
\end{definition}

\begin{lemma}[Zone length]
\lean{Complexity.payloadZone_length}
\label{Complexity.payloadZone_length}
\leanok
\uses{Complexity.OblPayload, Complexity.payloadZone}
\difficulty{3}
For every payload family, starting coordinate, caching flag, and length $n$, the zone
of $n$ rows is a list of length exactly $n$; every guide cell thus contributes exactly
one parallel data row.
\end{lemma}

\begin{lemma}[Full forward sweep]
\lean{Complexity.payloadForward_zone}
\label{Complexity.payloadForward_zone}
\leanok
\uses{Complexity.OblPayload, Complexity.payloadForward, Complexity.payloadForward_row, Complexity.payloadZone, Turing.FinTM.sweepFold}
\difficulty{4}
For every payload family, coordinate $z$, and length $n$, folding the forward
transducer with initial register the payloads at $z - 1$ over the uncached zone of $n$
rows starting at $z$ ends with register the payloads at $z + n - 1$ and produces the
cached zone of the same $n$ rows; all left-neighbor caches are installed.
\end{lemma}

\begin{lemma}[Full backward sweep]
\lean{Complexity.payloadBackward_zone}
\label{Complexity.payloadBackward_zone}
\leanok
\uses{Complexity.OblPayload, Complexity.payloadBackward, Complexity.payloadBackward_row, Complexity.payloadZone, Turing.FinTM.sweepFold, Turing.FinTM.sweepFold_append}
\difficulty{4}
For every payload family $f$, every vector of directions $d$, every coordinate $z$,
and every length $n$, folding the backward transducer with directions $d$ and initial
register the payloads at $z + n$ over the reversed cached zone of $n$ rows starting at
$z$ ends with register the payloads at $z$ and produces the reversed uncached zone of
the family shifted per tape by $d$; every row is shifted as prescribed.
\end{lemma}

\begin{lemma}[Two sweeps realize one source action]
\lean{Complexity.payload_sweeps_source}
\label{Complexity.payload_sweeps_source}
\leanok
\uses{Complexity.payloadBackward, Complexity.payloadBackward_zone, Complexity.payloadZone, Complexity.sourcePayload, Complexity.sourcePayload_apply, Complexity.sourceShift, Complexity.sourceWrittenPayload, Turing.Action, Turing.Action.apply, Turing.Cfg, Turing.FinTM.sweepFold}
\difficulty{3}
Let $c$ be a configuration of a binary machine, $\alpha$ an action, $z$ a coordinate,
and $n$ a width.  Folding the backward transducer with the shift directions of
$\alpha$, seeded with the written payloads of $c$ and $\alpha$ at $z + n$, over the
reversed cached zone of written payloads starting at $z$ yields the written payloads
at $z$ together with the reversed uncached zone of the head-centered payloads of the
configuration obtained by applying $\alpha$ to $c$; the two finite transductions thus
realize one source action on the whole fixed layout.
\end{lemma}

\begin{lemma}[Blank payloads outside the input]
\lean{Complexity.inputPayload_outside}
\label{Complexity.inputPayload_outside}
\leanok
\uses{Complexity.blankPayload, Complexity.inputPayload, Turing.FinTM.bufferTape}
\difficulty{4}
For every binary word $x$ and every integer coordinate $z$ with $z < -1$ or
$z > |x|$, the virtual input payload of $x$ at $z$ is the blank payload; outside both
virtual boundaries the head-centered input tape is blank.
\end{lemma}

\begin{lemma}[Support radius of the head-centered layout]
\lean{Complexity.sourcePayload_support}
\label{Complexity.sourcePayload_support}
\leanok
\uses{Complexity.blankPayload, Complexity.inputPayload_outside, Complexity.sourcePayload, Turing.FinTM, Turing.FinTM.source_bounds, Turing.MultiTapeTM.initCfg, Turing.MultiTapeTM.runFrom}
\difficulty{4}
Let $M$ be a binary machine, $x$ an input, and $t, B$ natural numbers with $t \le B$
and $|x| + 1 \le B$.  Then for every integer coordinate $z$ with $z < -3B$ or
$z > 3B$ and every virtual tape coordinate, the head-centered payload of the
configuration reached by running $M$ for $t$ steps from its initial configuration on
$x$ is the blank payload.  This bound uses the initialized run, not arbitrary starting
configurations.
\end{lemma}

\begin{lemma}[Last-output register under a one-symbol append]
\lean{Complexity.lastOutput_append}
\label{Complexity.lastOutput_append}
\leanok
\difficulty{2}
For every bit list $w$ and every optional bit $b$, the last entry, defaulting to
false, of $w$ extended by the zero- or one-element list determined by $b$ equals $b$
when $b$ is present, and the last entry of $w$ defaulting to false otherwise; a
one-bit last-output register can therefore be maintained in finite control.
\end{lemma}

\begin{lemma}[The answer sits in the last output bit at the padded budget]
\lean{Complexity.sourceAnswer_at_budget}
\label{Complexity.sourceAnswer_at_budget}
\leanok
\uses{Turing.FinTM, Turing.FinTM.ComputesInTime.mono, Turing.FinTM.DecidesInTime, Turing.FinTM.computesInTime_iff, Turing.MultiTapeTM.indicator, Turing.MultiTapeTM.initCfg, Turing.MultiTapeTM.runFrom}
\difficulty{3}
Let $M$ be a binary machine that decides a language $L$ over the binary alphabet
within time $a \cdot T(n)$ on inputs of length $n$.  Then for every input $x$, after
running $M$ for $(a + 1)(T(|x|) + 1)$ steps from its initial configuration on $x$, the
last bit of the output word, defaulting to false, equals the indicator of the
membership of $x$ in $L$.
\end{lemma}

\begin{definition}[Captured clock tape]
\lean{Complexity.clockTape}
\label{Complexity.clockTape}
\leanok
\uses{Complexity.OblSymbol, Turing.FinTM.bufferTape}
For a bit word $w$, the captured clock tape carries a permanent origin sentinel at
coordinate $0$ and, at each coordinate $z$ with $1 \le z \le |w|$, the bit symbol of
the corresponding bit of $w$; every other cell is blank.
\end{definition}

\begin{lemma}[Initial captured tape]
\lean{Complexity.clockTape_nil}
\label{Complexity.clockTape_nil}
\leanok
\uses{Complexity.clockTape}
\difficulty{2}
The captured clock tape of the empty word is blank everywhere except for the origin
sentinel at coordinate $0$.
\end{lemma}

\begin{lemma}[Capturing one emitted bit]
\lean{Complexity.clockTape_append}
\label{Complexity.clockTape_append}
\leanok
\uses{Complexity.clockTape, Turing.FinTM.bufferTape_append}
\difficulty{4}
For every bit word $w$ and every bit $b$, the captured clock tape of $w$ extended by
$b$ differs from the captured tape of $w$ in exactly one cell: the cell at coordinate
$|w| + 1$, which becomes the bit symbol of $b$.
\end{lemma}

\begin{definition}[Clock-stage configuration]
\lean{Complexity.clockStageCfg}
\label{Complexity.clockStageCfg}
\leanok
\uses{Complexity.OblPhase, Complexity.OblSymbol, Complexity.clockTape, Complexity.oblEmbed, Turing.Cfg, Turing.FinTM}
Let $W$ be a binary machine with $k$ work tapes, $a$ a padding constant, and $c$ a
configuration of $W$ on an input $x$.  The clock-stage configuration of the schedule
controller on the block-encoded input has control state the clock-phase encoding of
the state of $c$ (a designated reset phase once $c$ has halted), input head at the
position of the input head of $c$, the $k$ leading work tapes carrying the
bit-encoded work tapes of $c$, the budget tape capturing the entire append-only
output of $c$ with head one past the captured word, the unary tape holding only an
origin sentinel with head at coordinate one, a blank guide tape with head at the
origin, and empty output.
\end{definition}

\begin{lemma}[Masked input read of the clock stage]
\lean{Complexity.clockStageCfg_input}
\label{Complexity.clockStageCfg_input}
\leanok
\uses{Complexity.clockStageCfg, Turing.Cfg, Turing.Cfg.inputSymbol, Turing.FinTM}
\difficulty{3}
Let $W$ be a binary machine, $a$ a padding constant, and $c$ a configuration of $W$ on
an input $x$.  After forgetting the bit values (retaining only whether a symbol is
present), the symbol read by the input head of the clock-stage configuration of $c$
agrees with the correspondingly masked native read of $c$.
\end{lemma}

\begin{lemma}[Work reads of the clock stage]
\lean{Complexity.clockStageCfg_work}
\label{Complexity.clockStageCfg_work}
\leanok
\uses{Complexity.clockBit, Complexity.clockStageCfg, Turing.Cfg, Turing.Cfg.workTapeSymbols, Turing.FinTM}
\difficulty{3}
Let $W$ be a binary machine, $a$ a padding constant, and $c$ a configuration of $W$.
Decoding the bit symbols read on the $k$ leading work tapes of the clock-stage
configuration of $c$ recovers exactly the work-tape reads of $c$.
\end{lemma}

\begin{lemma}[One live clock step is one schedule step]
\lean{Complexity.clockStageCfg_step}
\label{Complexity.clockStageCfg_step}
\leanok
\uses{Complexity.clockStageCfg, Complexity.clockStageCfg_input, Complexity.clockStageCfg_work, Complexity.clockTape_append, Complexity.maskedClock, Complexity.obliviousSchedule, Turing.Action.apply, Turing.Cfg, Turing.Cfg.inputSymbol, Turing.Cfg.workTapeSymbols, Turing.FinTM, Turing.MultiTapeTM.step, Turing.moveInputPos}
\difficulty{5}
Let $W$ be a binary machine, $a$ a padding constant, and $c$ a live configuration of
$W$ on an input $x$ (its control state is some state $q$).  Then one step of the
schedule controller from the clock-stage configuration of $c$ equals the clock-stage
configuration of one step, from $c$, of the masked clock machine obtained from $W$ by
ignoring the input bit values; in particular any bit emitted by that step, including a
final emission, is captured on the budget tape, and administrative work begins only
after this transition.
\end{lemma}

\begin{lemma}[Exhaustion of the three administrative indices]
\lean{Complexity.finThree_cases}
\label{Complexity.finThree_cases}
\leanok
\difficulty{2}
Every element of the canonical three-element index type equals $0$, $1$, or $2$.
\end{lemma}

\begin{lemma}[Initialization of the clock stage]
\lean{Complexity.clockStageCfg_init}
\label{Complexity.clockStageCfg_init}
\leanok
\uses{Complexity.clockStageCfg, Complexity.clockTape_nil, Complexity.finThree_cases, Complexity.maskedClock, Complexity.oblAction, Complexity.oblEmbed, Complexity.obliviousSchedule, Turing.Action.apply, Turing.FinTM, Turing.MultiTapeTM.initCfg, Turing.MultiTapeTM.step}
\difficulty{4}
Let $W$ be a binary machine, $a$ a padding constant, and $x$ a binary input.  One step
of the schedule controller from its initial configuration on the block-encoded input
equals the clock-stage configuration of the initial configuration on $x$ of the masked
clock machine obtained from $W$; this single transition installs both sentinels and
starts the clock.
\end{lemma}

\begin{lemma}[Clock-stage run up to the first halt]
\lean{Complexity.clockStageCfg_run}
\label{Complexity.clockStageCfg_run}
\leanok
\uses{Complexity.clockStageCfg, Complexity.clockStageCfg_init, Complexity.clockStageCfg_step, Complexity.maskedClock, Complexity.oblEmbed, Complexity.obliviousSchedule, Turing.FinTM, Turing.MultiTapeTM.initCfg, Turing.MultiTapeTM.runFrom, Turing.MultiTapeTM.runFrom_succ_eq_step'}
\difficulty{4}
Let $W$ be a binary machine, $a$ a padding constant, $x$ a binary input, and $\tau$ a
time bound such that the masked clock machine obtained from $W$, started on $x$, is
still live at every time strictly before $\tau$.  Then for every $t \le \tau$, running
the schedule controller for $t + 1$ steps from its initial configuration on the
block-encoded input equals the clock-stage configuration of the masked clock run for
$t$ steps from its initial configuration on $x$; every emitted budget bit is thereby
represented.
\end{lemma}

\begin{lemma}[The schedule captures the budget word]
\lean{Complexity.clockStageCfg_captures}
\label{Complexity.clockStageCfg_captures}
\leanok
\uses{Complexity.clockStageCfg, Complexity.clockStageCfg_run, Complexity.guideTape, Complexity.maskedClock, Complexity.maskedClock_computes, Complexity.oblEmbed, Complexity.obliviousSchedule, Turing.Cfg, Turing.FinTM, Turing.FinTM.ComputesInTime, Turing.FinTM.ComputesInTime.output_unique, Turing.FinTM.computesInTime_iff, Turing.MultiTapeTM.initCfg, Turing.MultiTapeTM.runFrom}
\difficulty{5}
Let $W$ be a binary machine, let $a$ and $b$ be constants, and let a time bound $T$ be
such that on every input $x$, the machine $W$ outputs the binary representation of
$T(|x|)$ within $b\,(T(|x|)+1)$ steps.  Then for every input $x$ there exist a time
$\tau \le b\,(T(|x|)+1)$ and a halted configuration $c$ of $W$ on $x$ whose output is
the binary representation of $T(|x|)$, such that running the schedule controller for
$\tau + 1$ steps from its initial configuration on the block-encoded input equals the
clock-stage configuration of $c$.  No assumption that $W$ is oblivious is used.
\end{lemma}

\begin{definition}[Fixed sweep guide]
\lean{Complexity.guideTape}
\label{Complexity.guideTape}
\leanok
\uses{Complexity.OblSymbol}
For a radius $R$, the guide tape carries a left edge marker at coordinate $-R - 1$, a
right edge marker at coordinate $R + 1$, a marked origin at coordinate $0$, an
interior marker at every other coordinate $z$ with $-R \le z \le R$, and blanks
elsewhere.
\end{definition}

\begin{lemma}[Location of the left edge]
\lean{Complexity.guideTape_left}
\label{Complexity.guideTape_left}
\leanok
\uses{Complexity.guideTape}
\difficulty{2}
For every radius $R$ and every integer coordinate $z$, the guide tape of radius $R$
carries the left edge marker at $z$ if and only if $z = -R - 1$.
\end{lemma}

\begin{lemma}[Location of the right edge]
\lean{Complexity.guideTape_right}
\label{Complexity.guideTape_right}
\leanok
\uses{Complexity.guideTape}
\difficulty{2}
For every radius $R$ and every integer coordinate $z$, the guide tape of radius $R$
carries the right edge marker at $z$ if and only if $z = R + 1$.
\end{lemma}

\begin{lemma}[Location of the origin marker]
\lean{Complexity.guideTape_origin}
\label{Complexity.guideTape_origin}
\leanok
\uses{Complexity.guideTape}
\difficulty{2}
For every radius $R$ and every integer coordinate $z$, the guide tape of radius $R$
carries the origin marker at $z$ if and only if $z = 0$.
\end{lemma}

\begin{definition}[Macrostep configuration]
\lean{Complexity.macroCfg}
\label{Complexity.macroCfg}
\leanok
\uses{Complexity.OblPhase, Complexity.OblSymbol, Turing.Cfg, Turing.FinTM}
Let $W$ be a binary machine, $a$ a padding constant, and let a base configuration of
the schedule controller be given together with an optional control phase and integers
$u$ and $g$.  The macrostep configuration agrees with the base on every tape content,
on the input head, on the budget head, and on the output; only the control phase
(set to the given one), the unary-counter head (at $u$), and the guide head (at $g$)
vary.
\end{definition}

\begin{lemma}[Unary read of a macrostep configuration]
\lean{Complexity.macroCfg_unary}
\label{Complexity.macroCfg_unary}
\leanok
\uses{Complexity.OblPhase, Complexity.OblSymbol, Complexity.macroCfg, Turing.Cfg, Turing.Cfg.workTapeSymbols, Turing.FinTM}
\difficulty{1}
In every macrostep configuration with unary head at $u$, the symbol read on the unary
tape is the base configuration's unary tape content at coordinate $u$.
\end{lemma}

\begin{lemma}[Guide read of a macrostep configuration]
\lean{Complexity.macroCfg_guide}
\label{Complexity.macroCfg_guide}
\leanok
\uses{Complexity.OblPhase, Complexity.OblSymbol, Complexity.macroCfg, Turing.Cfg, Turing.Cfg.workTapeSymbols, Turing.FinTM}
\difficulty{1}
In every macrostep configuration with guide head at $g$, the symbol read on the guide
tape is the base configuration's guide tape content at coordinate $g$.
\end{lemma}

\begin{lemma}[Read-only actions move only the two macrostep heads]
\lean{Complexity.macroCfg_apply}
\label{Complexity.macroCfg_apply}
\leanok
\uses{Complexity.OblPhase, Complexity.OblSymbol, Complexity.finThree_cases, Complexity.macroCfg, Complexity.oblAction, Turing.Action.apply, Turing.Cfg, Turing.FinTM, Turing.moveInputPos_zero}
\difficulty{4}
Let a macrostep configuration with phase $q$, unary head $u$, and guide head $g$ be
given.  Applying the read-only schedule action that enters a phase $q'$, keeps the
input and budget heads fixed, writes nothing, and moves the unary and guide heads in
directions $d_u$ and $d_g$ yields the macrostep configuration with phase $q'$, unary
head $u + d_u$, and guide head $g + d_g$.
\end{lemma}

\begin{lemma}[Counter test]
\lean{Complexity.macroCfg_check}
\label{Complexity.macroCfg_check}
\leanok
\uses{Complexity.OblPhase, Complexity.OblSymbol, Complexity.macroCfg, Complexity.macroCfg_apply, Complexity.macroCfg_unary, Complexity.obliviousSchedule, Turing.Action.apply, Turing.Cfg, Turing.FinTM, Turing.MultiTapeTM.step}
\difficulty{3}
For every macrostep configuration in the counter-test phase with unary head $u$ and
guide head $g$: if the base's unary tape holds a unit marker at $u$, one step of the
schedule controller enters the leftward-seek phase; otherwise it halts.  In both
cases every head stays in place.
\end{lemma}

\begin{lemma}[Outward scan step]
\lean{Complexity.macroCfg_seek}
\label{Complexity.macroCfg_seek}
\leanok
\uses{Complexity.OblPhase, Complexity.OblSymbol, Complexity.guideTape, Complexity.guideTape_left, Complexity.macroCfg, Complexity.macroCfg_apply, Complexity.macroCfg_guide, Complexity.obliviousSchedule, Turing.Action.apply, Turing.Cfg, Turing.FinTM, Turing.MultiTapeTM.step}
\difficulty{3}
Suppose the base's guide tape is the fixed guide of radius $R$.  From the
leftward-seek phase with guide head at $g$, one step of the schedule controller turns
into the forward phase and moves the guide head right exactly when $g = -R - 1$, and
otherwise stays in the seek phase and moves the guide head left.
\end{lemma}

\begin{lemma}[Forward scan step]
\lean{Complexity.macroCfg_forward}
\label{Complexity.macroCfg_forward}
\leanok
\uses{Complexity.OblPhase, Complexity.OblSymbol, Complexity.guideTape, Complexity.guideTape_right, Complexity.macroCfg, Complexity.macroCfg_apply, Complexity.macroCfg_guide, Complexity.obliviousSchedule, Turing.Action.apply, Turing.Cfg, Turing.FinTM, Turing.MultiTapeTM.step}
\difficulty{3}
Suppose the base's guide tape is the fixed guide of radius $R$.  From the forward
phase with guide head at $g$, one step of the schedule controller turns into the
backward phase and moves the guide head left exactly when $g = R + 1$, and otherwise
stays in the forward phase and moves the guide head right.
\end{lemma}

\begin{lemma}[Backward scan step]
\lean{Complexity.macroCfg_backward}
\label{Complexity.macroCfg_backward}
\leanok
\uses{Complexity.OblPhase, Complexity.OblSymbol, Complexity.guideTape, Complexity.guideTape_left, Complexity.macroCfg, Complexity.macroCfg_apply, Complexity.macroCfg_guide, Complexity.obliviousSchedule, Turing.Action.apply, Turing.Cfg, Turing.FinTM, Turing.MultiTapeTM.step}
\difficulty{3}
Suppose the base's guide tape is the fixed guide of radius $R$.  From the backward
phase with guide head at $g$, one step of the schedule controller turns into the
return-to-center phase and moves the guide head right exactly when $g = -R - 1$, and
otherwise stays in the backward phase and moves the guide head left.
\end{lemma}

\begin{lemma}[Return scan step]
\lean{Complexity.macroCfg_center}
\label{Complexity.macroCfg_center}
\leanok
\uses{Complexity.OblPhase, Complexity.OblSymbol, Complexity.guideTape, Complexity.guideTape_origin, Complexity.macroCfg, Complexity.macroCfg_apply, Complexity.macroCfg_guide, Complexity.obliviousSchedule, Turing.Action.apply, Turing.Cfg, Turing.FinTM, Turing.MultiTapeTM.step}
\difficulty{3}
Suppose the base's guide tape is the fixed guide of radius $R$.  From the
return-to-center phase with guide head at $g$, one step of the schedule controller
commits exactly at the origin: when $g = 0$ it enters the counter-test phase and
advances the unary head by one, and otherwise it stays in the return phase and moves
the guide head right.
\end{lemma}

\begin{lemma}[Exact prefix of the outward scan]
\lean{Complexity.macroCfg_seek_run}
\label{Complexity.macroCfg_seek_run}
\leanok
\uses{Complexity.OblPhase, Complexity.OblSymbol, Complexity.guideTape, Complexity.macroCfg, Complexity.macroCfg_seek, Complexity.obliviousSchedule, Turing.Cfg, Turing.FinTM, Turing.MultiTapeTM.runFrom, Turing.MultiTapeTM.runFrom_succ_eq_step'}
\difficulty{4}
Suppose the base's guide tape is the fixed guide of radius $R$ and let $n \le R + 1$.
Running the schedule controller for $n$ steps from the leftward-seek phase with guide
head at the origin stays in the seek phase and moves the guide head to $-n$.
\end{lemma}

\begin{lemma}[Exact prefix of the forward scan]
\lean{Complexity.macroCfg_forward_run}
\label{Complexity.macroCfg_forward_run}
\leanok
\uses{Complexity.OblPhase, Complexity.OblSymbol, Complexity.guideTape, Complexity.macroCfg, Complexity.macroCfg_forward, Complexity.obliviousSchedule, Turing.Cfg, Turing.FinTM, Turing.MultiTapeTM.runFrom, Turing.MultiTapeTM.runFrom_succ_eq_step'}
\difficulty{4}
Suppose the base's guide tape is the fixed guide of radius $R$ and let
$n \le 2R + 1$.  Running the schedule controller for $n$ steps from the forward phase
with guide head at $-R$ stays in the forward phase and moves the guide head to
$-R + n$.
\end{lemma}

\begin{lemma}[Exact prefix of the backward scan]
\lean{Complexity.macroCfg_backward_run}
\label{Complexity.macroCfg_backward_run}
\leanok
\uses{Complexity.OblPhase, Complexity.OblSymbol, Complexity.guideTape, Complexity.macroCfg, Complexity.macroCfg_backward, Complexity.obliviousSchedule, Turing.Cfg, Turing.FinTM, Turing.MultiTapeTM.runFrom, Turing.MultiTapeTM.runFrom_succ_eq_step'}
\difficulty{4}
Suppose the base's guide tape is the fixed guide of radius $R$ and let
$n \le 2R + 1$.  Running the schedule controller for $n$ steps from the backward phase
with guide head at $R$ stays in the backward phase and moves the guide head to
$R - n$.
\end{lemma}

\begin{lemma}[Exact prefix of the return scan]
\lean{Complexity.macroCfg_center_run}
\label{Complexity.macroCfg_center_run}
\leanok
\uses{Complexity.OblPhase, Complexity.OblSymbol, Complexity.guideTape, Complexity.macroCfg, Complexity.macroCfg_center, Complexity.obliviousSchedule, Turing.Cfg, Turing.FinTM, Turing.MultiTapeTM.runFrom, Turing.MultiTapeTM.runFrom_succ_eq_step'}
\difficulty{4}
Suppose the base's guide tape is the fixed guide of radius $R$ and let $n \le R$.
Running the schedule controller for $n$ steps from the return-to-center phase with
guide head at $-R$ stays in the return phase and moves the guide head to $-R + n$.
\end{lemma}

\begin{lemma}[Exact macrostep duration]
\lean{Complexity.macroCfg_cycle}
\label{Complexity.macroCfg_cycle}
\leanok
\uses{Complexity.OblPhase, Complexity.OblSymbol, Complexity.guideTape, Complexity.macroCfg, Complexity.macroCfg_backward, Complexity.macroCfg_backward_run, Complexity.macroCfg_center, Complexity.macroCfg_center_run, Complexity.macroCfg_check, Complexity.macroCfg_forward, Complexity.macroCfg_forward_run, Complexity.macroCfg_seek, Complexity.macroCfg_seek_run, Complexity.obliviousSchedule, Turing.Cfg, Turing.FinTM, Turing.MultiTapeTM.runFrom, Turing.MultiTapeTM.runFrom_add, Turing.MultiTapeTM.runFrom_succ_eq_step'}
\difficulty{5}
Suppose the base's guide tape is the fixed guide of radius $R$ and its unary tape
holds a unit marker at coordinate $u$.  Then running the schedule controller for
exactly $6R + 8$ steps from the counter-test phase with unary head at $u$ and guide
head at the origin executes the entire fixed path of one macrostep and returns to the
counter-test phase with unary head at $u + 1$ and guide head at the origin.
\end{lemma}

\begin{lemma}[Repeated macrosteps over a unary segment]
\lean{Complexity.macroCfg_repeat}
\label{Complexity.macroCfg_repeat}
\leanok
\uses{Complexity.OblPhase, Complexity.OblSymbol, Complexity.guideTape, Complexity.macroCfg, Complexity.macroCfg_cycle, Complexity.obliviousSchedule, Turing.Cfg, Turing.FinTM, Turing.MultiTapeTM.runFrom, Turing.MultiTapeTM.runFrom_add}
\difficulty{4}
Suppose the base's guide tape is the fixed guide of radius $R$ and its unary tape
holds unit markers at all coordinates $1, 2, \dots, B$.  Then for every $j \le B$,
running the schedule controller for $j\,(6R + 8)$ steps from the counter-test phase
with unary head at $1$ and guide head at the origin reaches the counter-test phase
with unary head at $j + 1$ and guide head at the origin, regardless of all simulated
data and of source halting behavior.
\end{lemma}

\begin{lemma}[End of the fixed schedule]
\lean{Complexity.macroCfg_finish}
\label{Complexity.macroCfg_finish}
\leanok
\uses{Complexity.OblPhase, Complexity.OblSymbol, Complexity.guideTape, Complexity.macroCfg, Complexity.macroCfg_check, Complexity.macroCfg_repeat, Complexity.obliviousSchedule, Turing.Cfg, Turing.FinTM, Turing.MultiTapeTM.runFrom, Turing.MultiTapeTM.runFrom_succ_eq_step'}
\difficulty{3}
Suppose the base's guide tape is the fixed guide of radius $R$, its unary tape holds
unit markers at all coordinates $1, 2, \dots, B$, and the unary cell at coordinate
$B + 1$ is blank.  Then running the schedule controller for $B\,(6R + 8) + 1$ steps
from the counter-test phase with unary head at $1$ and guide head at the origin ends
in the halted phase with unary head at $B + 1$ and guide head at the origin: the
schedule finishes after exactly the budgeted macrosteps plus one final counter test.
\end{lemma}
